\chapter{Conclusion and Future Work}

\section{Summary}

This thesis studied learning-augmented algorithms for online bipartite
matching by first building a common experimental foundation and then
explaining what it revealed. On a single harness --- validated against
the published Borodin et al.~study (Chapter 3) --- we compared the
advice-free baselines, MinPredictedDegree and its augmentations, and the
test-and-fallback algorithms, across synthetic graphs, six real-world
graphs, and real request traces (Chapters 4--7).

Across every setting the same picture appeared: the advice-free baseline
is already near-optimal on average-case inputs, so the
\emph{consistency} upside of a good prediction is small; unguarded
prediction-following crashes \emph{below} the baseline under bad
predictions; and the value of the sophisticated algorithms is downside
protection, delivered by either a structural or an adaptive robustness
mechanism with distinct trade-offs.

Two sub-questions were sharpened along the way. The (small) loss is
governed by a Kendall-\(\tau\) order error rather than by magnitude,
order-dependence itself being a prior result of Aamand, Chen and Indyk
(Chapter 5). The adaptive test has a threshold-calibration pathology and
a resolution limit (Chapter 6). We also report two directions that did
not work: learning the predictor for extra performance, and finding a
serving regime where predictions genuinely help (Chapter 8).

The recurring wall raises an obvious question: is it an accident of the
inputs we chose, or is it forced? The next section gives our answer in
outlook form. The single sentence that unifies the thesis is: \textbf{on
average-case online matching, predictions are robustness insurance
rather than a performance lever --- and the upside they offer is smaller
than the price of finding out whether to trust them.}

\section{A theoretical outlook: the price of testing
advice}

This section does two things and nothing more: it proves one inequality,
and then reads our own experiments through it. Nothing beyond that
inequality is claimed as a theorem, and the reading that follows it is
an interpretation of measurements rather than a proof.

With the scope fixed: the experiments say that no \emph{practical}
acceptance threshold captures the upside safely (§6.3, Figure 6.4). The
question here is whether a decision rule of some \emph{other} kind
could, on the inputs where the wall stands.

\begin{quote}
\textbf{A trade-off inequality.} Let \(G\) and \(\mathrm{Bd}\) be two
instance distributions sharing the \emph{same} advice, such that
following the advice gains \(\delta\) under \(G\) and loses \(\Delta\)
under \(\mathrm{Bd}\) relative to the advice-free baseline, and let
\(\gamma_k\) be the total-variation distance between the laws of their
length-\(k\) prefixes --- informally, the best accuracy any test can
reach when it sees only the first \(k\) arrivals. Then every
test-and-fallback algorithm --- deciding by \emph{any} measurable rule
on its prefix --- satisfies \[(1-\eta_c)\;\le\;\eta_r+\gamma_k+o(1),\]
where \(\eta_c\) is the fraction of the upside it forgoes under \(G\)
and \(\eta_r\) its robustness loss under \(\mathrm{Bd}\) as a fraction
of \(\Delta\).
\end{quote}

The proof is a short conditioning argument: capturing the upside forces
the algorithm to follow with high probability under \(G\); robustness
forces fallback under \(\mathrm{Bd}\); and no function of the prefix can
behave differently on two prefix distributions that are statistically
\(\gamma_k\)-close. The inequality thus converts ``consistent \emph{and}
robust'' into a question about \emph{sample complexity}: how long must
the prefix be before it distinguishes advice worth following from advice
worth rejecting?

The reading applies to the rare-resource instances that produce Figure
6.4 --- a \emph{decomposable} family, in which the instance splits into
independent cells and the advice acts on each of them separately --- and
it is a \textbf{budget--stakes} relation. Write \(\delta\) for the
\emph{stakes}: what good advice gains over the baseline. In the
companion development, which is not proved here, the decision costs a
prefix of \(k^* = \tilde\Theta(\theta/\delta^2)\) --- the inverse square
of the stakes, scaled by a contention parameter \(\theta\) that on these
instances is of the same order as the baseline slack
\(1-\rho_{\mathrm{base}}\) --- and that budget is met by a simple
\emph{directional} statistic (do the prefix arrivals agree with the
advice's predictions more often than they contradict them?). The stakes
are in turn capped by the same slack,
\(\delta = O(1-\rho_{\mathrm{base}})\): advice can only win what the
baseline leaves on the table. Substituting, the reading predicts that on
strong-baseline instances the required prefix exceeds the whole instance
--- that an upside below
\(\Theta\bigl(\sqrt{(1-\rho_{\mathrm{base}})/n}\bigr)\) would not be
captured by any rule at any prefix length \(k \le n\). At the parameters
of Chapter 4 (\(\rho_{\mathrm{base}}\approx0.99\), \(n=2000\)) that
threshold is \(\approx0.004\), and the upsides we measured are of
exactly this order (F3): the empirical wall sits where the reading says
it should. Whether the same relation holds beyond decomposable families
is open (§10.4). This is Figure 6.4 read quantitatively: the potential
upside and the capturable upside separate as the baseline strengthens,
because the stakes shrink faster than the test's resolution improves.

Two honest notes bound the scope of this outlook. First, the budget law
cuts both ways: where the stakes are large (a weak baseline), the
directional statistic is cheap and following good advice is easy --- the
wall is a statement about strong-baseline, average-case inputs, not
about testing in general. In particular the wall is not a consequence of
distribution testing being expensive --- a directional statistic is
cheap --- but of the stakes being small; the \(\ell_1\)-threshold
blindness of §6.3 comes from testing the \emph{distance} rather than the
\emph{payoff}. Second, this thesis claims only the inequality displayed
above and the quantitative reading of its own experiments; the two-sided
law and its exact scope are not established here (§10.3).

\section{Limitations}

\begin{itemize}
\tightlist
\item
  \textbf{Input model.} We work in the known-i.i.d. model. Because every
  known-i.i.d. instance is also a random-order instance, guarantees
  proved in the random-order model carry over to ours (but not
  conversely, §2.1); the empirical wall, however, is an
  \emph{average-case} statement; we do not claim it for adversarial
  arrival order.
\item
  \textbf{Test model.} Following the original authors, the
  test-and-fallback experiments use an empirical-\(\ell_1\) surrogate
  for the (unimplemented) distribution tester; §10.2 explains why the
  surrogate's blindness is structural rather than an implementation
  artifact.
\item
  \textbf{Prediction-object heterogeneity.} The degree- and
  histogram-prediction families do not map onto every graph, which is
  why they are reported in parallel panels rather than one table.
\item
  \textbf{Data breadth.} Each real modality is exercised by one trace.
\item
  \textbf{Objective.} We evaluate matching size (goodput) only. On
  objectives where the advice-free baseline is \emph{not} near-optimal
  --- tail latency, per-type fairness, migration or recompute cost ---
  the picture could differ; our one probe in that direction (§8.2)
  closed the simplest such attempt but not the space (§10.4).
\item
  \textbf{Theory scope.} The thesis deliberately confines its theory to
  the outlook of §10.2 --- one proved trade-off inequality and a
  quantitative reading of the experiments. The full budget--stakes law
  is companion work in preparation, and no theorem beyond the stated
  inequality is claimed here.
\end{itemize}

\section{Future work}

The thesis brackets, rather than resolves, the settings in which
predictions might genuinely help online matching, and these are the
natural next directions.

\begin{itemize}
\tightlist
\item
  \textbf{Beyond average-case inputs.} The wall is an average-case
  phenomenon. Adversarial or non-stationary arrival orders, where the
  advice-free baseline is provably far from optimal, are where
  predictions should carry real value --- and where a \emph{positive}
  counterpart to this thesis's wall might be proved.
\item
  \textbf{Beyond throughput.} Objectives on which the baseline is not
  near-optimal --- tail latency, per-type fairness, migration or
  recompute cost --- may admit genuine with-predictions gains that the
  goodput objective forecloses; our serving SLO probe (Chapter 8) closed
  the simplest such attempt but not the space.
\item
  \textbf{Completing the theory.} Finishing the companion budget--stakes
  development --- the sharp two-sided law, the directional statistic as
  its matching upper bound, and its exact scope (decomposable families;
  whether a non-decomposable family can push the budget higher is open)
  --- would turn the outlook of §10.2 into a tight characterization.
\item
  \textbf{Better tests, honestly.} Whether a super-linear or amortized
  test, or a test that reuses online decisions as samples, can beat the
  stakes-squared budget of §10.2 is an open and practically motivated
  question.
\end{itemize}

\section{Closing}

Predictions are a powerful tool for online algorithms, but this thesis
is a study of their \emph{limits} on one well-understood problem. On
average-case online matching, the honest verdict has three parts. A
cheap, order-faithful predictor already captures nearly all there is to
capture. The sophisticated machinery earns its keep as insurance rather
than as performance. And finding out whether to trust a prediction costs
more, on these inputs, than the prediction is worth. Recognizing where
predictions cannot help is, we hope, as useful as knowing where they
can.
